%\usepackage[letterpaper,showframe]{geometry}	
% usuwanie nagłówków/stopek z\,dokumentu
\usepackage[skip=1pt,indent=0pt]{parskip} % ustawianie paragrafów
%\setstretch{1.5} % półtorej linii odległości między wierszami
\usepackage{titlesec} % ustawianie nagłówków
%\usepackage{fullpage} % redukcja marginesów
\usepackage{setspace} % odstępy między liniami
\usepackage{enumitem}% listy numerowane
%\usepackage{soul}% dodaje \ul
%\usepackage{ulem}% dodaje \uline
\usepackage{hanging}%obsługa akapitów
% --------------------------------------------------------------------------------------
% ------------------ USTAWIENIE NUMERÓW STRON ----------------------------------
\usepackage{fancyhdr}
\pagestyle{fancy}
\fancypagestyle{plain}{%
	\fancyhf{}% wyczyść wszystkie nagłówki i stopki
	\fancyfoot[LE,RO]{\footnotesize\thepage}% pozycja numeru strony LE - left on even, RO - right on odd
	\renewcommand{\headrulewidth}{0pt}%
	\renewcommand{\footrulewidth}{0pt}%
	}
%\renewcommand*{\thepage}{\footnotesize{\arabic{page}}}% mniejsza czcionka numeru strony, dodanie czegokolwiek przed numerem: {A--\arabic{page}}

% --------------------------------------------------------------------------------------
\newcommand{\margintext}[1]{\textit{\small #1}}
% przesunięcia poziome
\newcommand{\smallwidth}{\hskip 1em}
\newcommand{\medwidth}{\hskip 3pt}
\newcommand{\largewidth}{\hskip 12pt}
% przesunięcia pionowe
\newcommand{\minispace}{\vskip 4pt}
\newcommand{\linespace}{\vskip 7pt}
\newcommand{\paragraphspace}{\vskip 8pt}
\newcommand{\segmentspace}{\vskip 18pt}
\newcommand{\sectionspace}{\vskip 28pt}
% -----------------------------------------------------------------------------------
% --------------------------------- GREGORIO --------------------------------------
% -----------------------------------------------------------------------------------
\usepackage{gregoriotex}
\gresetlinecolor{gregoriocolor}% Make staff lines red; remove for black:
%\gresetgregoriofont{granapadano} %{gregorio}
\gresetheadercapture{annotation}{greannotation}{string}% odczytuje adnotacje nad inicjałem jako ciąg, co pozwala na jego renderowanie
\gresetheadercapture{commentary}{grecommentary}{string}% odczytuje komentarz jako ciąg, co pozwala na jego renderowanie; wstawianie komentarza: \grecommentary{}
\grechangedim{commentaryraise}{0.3cm}{scalable} % podnosi komentarz wystarczająco ponad górną linię
\grechangestyle{annotation}{\color{black}\small}% format adnotacji
\grechangestyle{commentary}{\color{gregoriocolor}\small\itshape}% format komentarza
\grechangestyle{abovelinestext}{\color{gregoriocolor}\small}% format napisu nad liniami
% --- inicjał
\grechangestyle{initial}{\greIni\fontsize{50}{36}\selectfont} % ustawia czcionkę zdefiniowaną w sekcji czcionek o zmniejsza wielkość inicjału do rozmiaru 46pt a baseline do 36pt
% -----
\grechangedim{intersyllablespacenotes}{0.05cm}{scalable}% Minimalna odległość między nutami różnych sylab.
% Set the space around the initial:
% See http://gregorio-project.github.io/gregoriotex/details.html for more details and options
\grechangedim{beforeinitialshift}{0.5mm}{scalable}
\grechangedim{afterinitialshift}{1.0mm}{scalable}
\grechangestaffsize{25}%rozmiar wyniku GregorioTEX

% umieszczenie zawartości składu ich unikalnego argumentu po lewej i prawej stronie każdej linii
\let\grelocalleftbox\localleftbox
\let\grelocalrightbox\localrightbox

%\newcommand{\separator}{\greseparator{3}{40}}
%\newfontface\vectis{VectisMonumentalW01-Regular}
%\def\greinitialformat#1{\fontsize{20}{20}\vectis #1}
% ------------------

\usepackage[hyperindex=true]{hyperref}

% --- SPIS TREŚCI -------
%\usepackage{imakeidx}
%\makeindex[title=Alfabetyczny spis program=texindy,options=-M styl.xdy]
%\usepackage{filecontents}
%\begin{filecontents*}{styl.xdy}
%;;; xindy style file
%(markup-locclass-list :open "\dotfill" :sep "")
%(define-attributes (( "textbf" "default" )) )
%(markup-locref   :attr  "textbf"     :open "\textbf{\hyperpage{" :close "}}")
%(markup-locref   :attr  "textit"     :open "\textit{\hyperpage{" :close "}}")
%(markup-locref   :attr  "textttt"     :open "\textttt{\hyperpage{" :close "}}")
%(markup-locref   :attr  "texttsc"     :open "\texttsc{\hyperpage{" :close "}}")
%(markup-locref   :attr  "default"     :open "\hyperpage{" :close "}")
%\end{filecontents*}
% -------------------------
%% --- UKŁAD DWUKOLUMNOWY Z TŁUMACZENIEM ---
%\usepackage{tcolorbox}
%\newtcolorbox{nutybox}{size=minimal, colback=white, colframe=red, arc=0mm, width=\linewidth, halign=left, valign=center}
%\newtcolorbox{tlumbox}{size=minimal, colback=white, colframe=red, arc=0mm, width=\linewidth, halign=left, valign=top, fontupper=\small\itshape} 
%\newcommand{\nuty}[1]{\begin{nutybox}#1 \end{nutybox}}
%\newcommand{\tlum}[1]{\begin{tlumbox}\tlumacz{#1}\end{tlumbox}}
%% ------------------------------------------------------
\usepackage{Tabbing}
\usepackage{longtable}
\usepackage[savepos]{zref}
% -----------------------------------------------------------------------------------
% ---------------------- CZCIONKI DOKUMENTU ------------------------
% -----------------------------------------------------------------------------------
\usepackage{fontspec}
\defaultfontfeatures{Ligatures=TeX}
\usepackage{polyglossia}
%\usepackage{ecclesiastic}% This package works only with pdfLaTeX 
\usepackage[autostyle=true]{csquotes}%generuje odpowiednie cudzysłowy dla danego języka
\usepackage{libertine}
\setmainfont[
  BoldFont          = LinLibertine_RZ.otf,
  ItalicFont        = LinLibertine_RI.otf,
  BoldItalicFont    = LinLibertine_RZI.otf,
  Ligatures=TeX]{LinLibertine_R.otf}
\setsansfont[
  BoldFont          = LinBiolinum_RB.otf,
  ItalicFont         = LinBiolinum_RI.otf,
  BoldItalicFont    = LinBiolinum_RBO.otf,
  Ligatures=TeX]{LinBiolinum_R.otf}
\setmonofont[
  BoldFont          = LinLibertine_MB.otf,
  ItalicFont        = LinLibertine_MO.otf,
  BoldItalicFont    = LinLibertine_MBO.otf,
  Ligatures=TeX]{LinLibertine_M.otf} 
\newfontface\greIni{Zallman Caps}% czcionka inicjału w czterolinii
%\newfontface\greIni{Preciosa}% czcionka inicjału w czterolinii
%\newfontface\greIni{Florana}% czcionka inicjału w czterolinii
\setdefaultlanguage[variant=ecclesiastic, hyphenation=liturgical,babelshorthands = true]{latin}
\setotherlanguage{polish}

%\setmainfont{Crimson}%ok
%\setmainfont{Old Standard TT}%ok
%\setmainfont{Crimson Pro}%ok
%\setmainfont{Berylium}%ok
%\setmainfont{Cardo}%ok
%\setmainfont{Antykwa Torunska}%ok
%\setmainfont{Alfios}%ok
%\setmainfont{MinionPro}
%\setmainfont{BaskervilleF}%ok
%\setmainfont{EB Garamond}%ok
% definiowanie fontów w  LuaLATEX':
%\setsansfont{CMU Sans Serif}
%\setmonofont{CMU Typewriter Text}
% ----
%\setmainfont{Libertinus Serif}%ok
%\setmainfont{TeX Gyre Termes}%ok - TimesNewRoman alternative
%\setmainfont{Calluna}
% ----
%\newfontfamily{\sekfont}{Times New Roman}
% -----------------------------------------------------------------------------------
% ------------------ CZCIONKA SEKCJI I SUBSEKCJI --------------------
\newcommand\sekcjafont{\fontsize{20}{16}\selectfont\scshape}
\newcommand\subsekcjafont{\fontsize{16}{10}\selectfont\scshape}
% -----------------------------------------------------------------------------------
% ------------------ CZCIONKA STRONY TYTUŁOWej ------------------
\newfontfamily{\tytfont}{JSL Ancient}
\newcommand\tytulfont[1]{\fontsize{36}{30}\selectfont\tytfont#1\par}
\newcommand\podtytulfont[1]{\fontsize{18}{10}\selectfont\tytfont#1\par}

% -----------------------------------------------------------------------------------
% ---------------------------- GRAFIKA I KOLORY---------------------------
% -----------------------------------------------------------------------------------
\usepackage{xcolor} % kolory
\definecolor{myRed}{RGB}{229,53,44}%
\usepackage{graphicx} % umożliwia import obrazów
\usepackage{pgfornament}% ornamenty graficzne
% -----------------------------------------------------------------------------------
% ----------------------------------- RUBRYKA---------------------------------
% -----------------------------------------------------------------------------------
\newcommand{\rubryka}[1]{
	\begin{spacing}{0.8}
	\begin{center}
		\parbox{0.9\textwidth}{\centering\small\textit{\red{#1}}}
	\end{center}
	\end{spacing}
	\linespace}
% -----------------------------------------------------------------------------------
% ---------------------------- STYLE NAGŁÓWKÓW------------------------
% -----------------------------------------------------------------------------------
% ----------------------------------- ROZDZIAŁ --------------------------------
% -----------------------------------------------------------------------------------
\titlespacing*{\chapter}{0pt}{-40pt}{0pt}	
\titleformat{\chapter}[hang]{\filcenter\scshape\bfseries\Large\sekcjafont}{}{0pt}{} % format nagłówków
\newcommand{\rozdzial}[2]{
	\begin{latin}
		\chapter{#1}
		\vspace{-1em}
		\rubryka{#2}
	\end{latin}}
% -----------------------------------------------------------------------------------
% -----------------------------------SEKCJA -------------------------------------
% -----------------------------------------------------------------------------------
\titlespacing*{\section}{0pt}{0pt}{0pt}	
\titleformat{\section}[hang]{\filcenter\scshape\normalfont\sekcjafont}{}{0pt}{} % format nagłówków
\newcommand{\sekcja}[2]{
	\begin{spacing}{0.6}
	\begin{center}
	\begin{latin}
		\segmentspace
		\section*{#1}
		\begin{minipage}[c]{0.88\linewidth}
			\rubryka{#2}
			%\centering\red{\small\textit{#2}}
		\end{minipage}
	\end{latin}
	\end{center}	
	\end{spacing}
	\linespace}	
% -----------------------------------------------------------------------------------
% -----------------------------------SUBSEKCJA -------------------------------
% -----------------------------------------------------------------------------------	
\titlespacing*{\subsection}{0pt}{0pt}{0pt}	
\titleformat{\subsection}[hang]{\scshape\normalfont\subsekcjafont}{}{0pt}{} % format podnagłówków
\newcommand{\subsekcja}[1]{
	\segmentspace
	\begin{latin}
		\subsection*{\red{#1}}
	\end{latin}		
	\segmentspace}
% -----------------------------------------------------------------------------------
% ----------------------------------- TYTUŁ -------------------------------------
% -----------------------------------------------------------------------------------
\newcommand{\tytul}[1]{
	\begin{latin}
		\red{\large\scshape{#1}}
		\linespace
	\end{latin}}
% -----------------------------------------------------------------------------------
% --------------------------------- PODTYTUŁ ---------------------------------
% -----------------------------------------------------------------------------------
\newcommand{\podtytul}[1]{
	\begin{latin}
		\red{\normalfont\scshape{#1}}
		\linespace
	\end{latin}}
% -----------------------------------------------------------------------------------
% ----------------------------------- DEKOR ------------------------------------
% -----------------------------------------------------------------------------------
\newcommand{\dekor}[2]{
	\begin{center}
		\vspace*{\fill}
		\includegraphics[width=#1\textwidth]{#2}
		\vspace*{\fill}
	\end{center}}
% -----------------------------------------------------------------------------------
% ---------------------------------- WERSETY ----------------------------------
% -----------------------------------------------------------------------------------
%\newcommand{\firstchar}[1]{\directlua{tex.sprint(unicode.utf8.sub('#1',1,1))}}%makro pobiera pierwszy znak ze stringu
\newcommand{\red}[1]{\textcolor{myRed}{#1}}
\newcommand{\mred}[1]{\leftskip7pt\makebox[3pt][r]{\textcolor{myRed}{\footnotesize{#1}}}}
\newcommand{\black}[1]{\textcolor{black}{#1}}
%\newcommand{\numer}[1]{\noindent\makebox[1em][l]{\textcolor{myRed}{#1}}}
% --------------------------------
\newcommand{\tlumacz}[1]{\begin{spacing}{0.5}\textpolish{\footnotesize\textit{#1}}\end{spacing}}% normalna czcionka justowana do lewej
\newcommand{\tlumaczR}[1]{\begin{spacing}{0.5}\mbox{}\hfill\textpolish{\footnotesize\textit{#1}}\end{spacing}}% mała czcionka justowana do prawej 
\newcommand{\inicjal}[2]{\textcolor{myRed}{\Large\textbf{#1}}{#2}}
\newcounter{licznik}
\newcommand{\licznikZero}{\setcounter{licznik}{0}}% zeruj licznik
\newcommand{\licznikUp}{\refstepcounter{licznik}}% zwiększ licznik 
\licznikZero
% -----------------------------------------------------------------------------------
% --------------------------------- KOLUMNY ---------------------------------
% -----------------------------------------------------------------------------------
\usepackage{paracol}
\wstawplik{conparlua}%---- usun konflikt między paracol i luacolor
\columnratio{0.5}
%\setcolumnwidth{.49\textwidth, .49\textwidth}
\setlength{\columnseprule}{0.4pt}
\setlength{\columnsep}{.02\textwidth}
\colseprulecolor{myRed}% wymagany pakiet xcolor
\footnotelayout{m}
 %------------------ kolorowe linie dzielące kolumny ------------------
\makeatletter 
\patchcmd\pcol@hfil
 {\copy\pcol@tempboxa}
 {\luacolorresetboxcolor\pcol@tempboxa\copy\pcol@tempboxa}{}{\fail}
\makeatother
 %-----------------------------------------------------------------------------------------------------
\newcommand{\BeginParallel}{\begin{spacing}{0.8}\begin{paracol}{2}\licznikZero}
\newcommand{\lat}[1]{\begin{latin}\ensurevspace{2\baselineskip}\begin{minipage}[t]{\linewidth}\ifnum \value{licznik}>0 \licznikUp \thelicznik .~#1 \else \licznikUp#1 \fi \end{minipage}\end{latin}\switchcolumn}
\newcommand{\latN}[1]{\begin{latin}\ensurevspace{2\baselineskip}\begin{minipage}[t]{\linewidth}#1\end{minipage}\end{latin}\switchcolumn}% kolumna bez numeracji wersów
\newcommand{\pol}[1]{\begin{polish}\begin{minipage}[t]{\linewidth}\textit{#1}\end{minipage}\end{polish}\switchcolumn*}
\newcommand{\polPs}[1]{\begin{polish}\begin{minipage}[t]{\linewidth}\footnotesize\textit{#1}\end{minipage}\end{polish}\switchcolumn*}% wersja dla tłumaczenia psalmów z mniejszą czcionką
\newcommand{\EndParallel}{\end{paracol}\end{spacing}}
% ------------------------------------------------------------------------------------------

% --- ZNAKI SPECJALNE ---
% 
%\newcommand †
% --------------------------
